\clearpage
\Article{362}
Mit diesen Grundlagen können wir die Primteiler der Funktion $x^\nu-1$ leicht bestimmen. Wir nehmen an, dass alle Teiler, die auch irgendeine Funktion
\SourceLeaf{241}
$x^\nu-1$ teilen, wobei $\nu'<\nu$, bereits gefunden sind, und wollen die übrigen untersuchen. Diese lassen sich alle im Ausdruck $(P,\rho^k)$ erfassen, wenn $P$ einer von ihnen ist und für $k$ alle Zahlen eingesetzt werden, die kleiner als $\nu$ und dazu teilerfremd sind.

Im Kapitel \textsc{vi} haben wir gezeigt, wie sich die primitiven Wurzeln der Gleichung $x^\nu=1$ so in Klassen aufteilen lassen, dass dieselbe Klasseneinteilung entsteht, wenn man alle durch Potenzen einer von ihnen ausdrückt, gleichgültig, welche primitive Wurzel man als Basis nimmt. Solche Zusammenstellungen von Wurzeln haben wir \emph{Perioden} genannt. Offenbar lassen sich die Funktionen $x,x^\alpha,x^\beta,x^\gamma$ usw., wobei $\alpha,\beta,\gamma$ usw. alle zu $\nu$ teilerfremden Zahlen bezeichnen, auf dieselbe Weise in Perioden zerlegen und jede größere Periode wiederum in kleinere, bis man schließlich zu Perioden der Form $x^k,x^{kp},x^{kpp}\dotstwo x^{kp^{m-1}}$ gelangt. Ist dies geschehen, so ist Folgendes klar:

$1$. Da jede Periode aus kleinsten Perioden der Form $x^k+x^{kp}+\text{usw.}$ zusammengesetzt ist, ergibt ihre Division durch irgendeine Primfunktion vom Grad $m$ einen Rest, der eine Zahl ist.

$2$. Alle Terme einer Periode lassen sich stets auf die Form $x^{\kappa\pdot a^\alpha b^\beta c^\gamma}\dotstwo$ bringen, wobei $\kappa.\ a,b,c\dotstwo$ feste Zahlen sind, während für $\alpha,\beta,\gamma\dotstwo$ alle Werte eingesetzt werden können. Daher geht die Periode in sich selbst über, wenn man $x$ durch $x^k$ ersetzt und $k$ die Form $a^\alpha b^\beta c^\gamma\dotstwo(\mathrm{mod.}\,\nu)$ hat. Hieraus erkennt man leicht, dass alle Funktionen $P,(P\;\rho^k)$ usw., wobei $k$ eine solche Zahl bezeichnet, bei der Division der Periode durch sie denselben Rest ergeben.

$3$. Nach Abzug eines solchen Restes lässt sich die Periode somit durch das Produkt aus allen Funktionen $(P,\rho^k)$ teilen.
